% Solution: Mas-Colell, Whinston & Green (1995) Exercise 8.D.2. Input by \solitem.
% New in the 50-item list. Check: check_A2.py.
\begin{exsol}{MWG Exercise 8.D.2}{263}{Advanced}{4}
\asked Prove that if exactly one strategy profile survives iterated removal of strictly dominated
strategies, that profile is a Nash equilibrium.

\method Iterated deletion of strictly dominated strategies (MWG \S8.B, pp.~236--241; Lecture~1,
slide~29) and Nash equilibrium (MWG Def.~8.D.1, pp.~246--247; Lecture~2, slide~14). The proof is
the slides' Proposition~2 (Lecture~2, slide~17): follow a chain of dominating strategies, along
which the payoff against the survivor's profile strictly rises, until it reaches the surviving
strategy. The game is finite (MWG \S8.B). The same chain works whether the dominators are pure
strategies (MWG Def.~8.B.2) or mixed (Def.~8.B.4); both versions are proved.

\solution
\textbf{Setting.} $\Gamma = [I, \{S_i\}, \{u_i\}]$, every $S_i$ finite. $S^0 = S$, and $S^{k+1}$
removes from $S^k$ every strategy strictly dominated given $S^k$ (by a pure strategy in $S^k_i$,
or by a mixed strategy $\sigma_i \in \Delta(S^k_i)$). The process stops after finitely many rounds
at $S^\infty$. Hypothesis: $S^\infty = \{s^*\}$.

\textbf{Claim.} For every player $i$ and every $s_i \in S_i$, $u_i(s_i^*, s_{-i}^*) \ge
u_i(s_i, s_{-i}^*)$.

\textbf{Step 1: one link of the chain.} Let $s_i \ne s_i^*$. Then $s_i \notin S^\infty_i$, so
$s_i$ was deleted in some round, say round $k+1$: there is a dominator, pure $t_i \in S^k_i$ or mixed
$\sigma_i \in \Delta(S^k_i)$, with
\[
  u_i(\sigma_i, s_{-i}) > u_i(s_i, s_{-i}) \quad \text{for every } s_{-i} \in S^k_{-i}. \tag{1}
\]
Since $s^*$ is never deleted, $s^*_{-i} \in S^k_{-i}$, and (1) holds at $s_{-i} = s^*_{-i}$:
\[
  \sum_{t_i \in S^k_i} \sigma_i(t_i)\, u_i(t_i, s^*_{-i}) = u_i(\sigma_i, s^*_{-i}) > u_i(s_i, s^*_{-i}).
\]
A weighted average exceeds $u_i(s_i, s^*_{-i})$ only if at least one of its terms does, so there is
a pure $t_i$ with $\sigma_i(t_i) > 0$ (in the pure case, $t_i$ itself) such that
\[
  t_i \in S^k_i \quad \text{and} \quad u_i(t_i, s^*_{-i}) > u_i(s_i, s^*_{-i}). \tag{2}
\]
Call $t_i$ a \emph{successor} of $s_i$. Since $t_i \in S^k_i$, it survived at least until the round
in which $s_i$ was deleted.

\textbf{Step 2: the chain.} Start at any $s_i^{(0)} = s_i \ne s_i^*$. While the current strategy
is not $s_i^*$, it was deleted at some round, so by Step~1 it has a successor; let $s_i^{(m+1)}$
be a successor of $s_i^{(m)}$. By (2),
\[
  u_i(s_i^{(0)}, s^*_{-i}) < u_i(s_i^{(1)}, s^*_{-i}) < u_i(s_i^{(2)}, s^*_{-i}) < \cdots
\]
The payoffs strictly increase, so no strategy appears twice; $S_i$ is finite, so the chain cannot go
on forever. It stops only at a strategy with no successor, that is, at a strategy that is never
deleted: $s_i^*$, the only one in $S^\infty_i$. Hence
\[
  u_i(s_i, s^*_{-i}) < u_i(s_i^*, s^*_{-i}) \quad \text{for every } s_i \ne s_i^*.
\]

\textbf{Step 3: conclusion.} No player gains by a unilateral pure deviation from $s^*$, so $s^*$ is
a Nash equilibrium (in fact a \emph{strict} one). Mixed deviations do not help either, since
$u_i(\sigma_i, s^*_{-i})$ is a weighted average of the pure payoffs, each at most
$u_i(s^*)$. $\blacksquare$

\textbf{Remarks.}
\begin{itemize}
  \item \emph{Uniqueness.} Every Nash equilibrium survives IESDS (Lecture~2, Proposition~1,
  slide~16), so $s^*$ is also the only pure Nash equilibrium; and in any mixed equilibrium every
  strategy in the support survives (MWG manual's argument below), so $s^*$ is the only Nash
  equilibrium of any kind.
  \item \emph{Finiteness} is used in Step~2: in an infinite game the chain of successors need not
  end at $s^*$, and a unique survivor need not be an equilibrium. (Outside the course books:
  Dufwenberg and Stegeman, \emph{Econometrica} 2002, study this.) Infinite games in the course where
  IESDS has a unique survivor, such as the Cournot duopoly of MWG 8.B.5, satisfy the conclusion by
  direct computation.
\end{itemize}

\interp If common knowledge of rationality alone pins down a single profile, that profile is
self-enforcing: no player can do better against what the others are predicted to do. Nash
equilibrium is weaker than IESDS in what it assumes about beliefs, and this result says the two
agree whenever IESDS is decisive, as in the campaigning game (Tadelis 4.7) or the discrete auction
(Tadelis 4.3).

\verify \texttt{check\_A2.py} generates random finite games (two players, two to four strategies,
integer payoffs) until it has 1{,}000 with a unique IESDS survivor under pure dominance, and checks
that each survivor is a strict Nash equilibrium and the only pure one; it also builds the chain of
Step~2 explicitly in each game and checks that payoffs rise along it and that it ends at $s^*_i$.
\emph{Comparison with the MWG solutions manual} (p.~8-9): the manual proves the converse
inclusion, that every (mixed) Nash equilibrium survives iterated strict dominance, and concludes
that the equilibrium must be the unique survivor. That route needs an equilibrium to exist in the
first place, which in a finite game is Nash's theorem (MWG Prop.~8.D.2) but is not stated in the
manual's answer. The direct chain argument above needs no existence theorem. The conclusions
agree.

\mistakes
\begin{itemize}
  \item Proving the converse (every Nash equilibrium survives IESDS) and stopping. That shows the
  survivor is the only \emph{candidate}; existence of an equilibrium is a separate step.
  \item Assuming the strategy that dominated $s_i$ survives to the end. It may itself be deleted
  later; the chain handles this.
  \item Comparing payoffs against all $s_{-i}$ instead of against $s^*_{-i}$. Only the inequality
  at $s^*_{-i}$ is needed, and it holds because $s^*_{-i}$ is never deleted.
\end{itemize}

\answer{Let $s^*$ be the unique survivor and $s_i \ne s_i^*$. When $s_i$ was deleted, its
dominator (or a pure strategy in its support) $t_i$ was still present and, since $s^*_{-i}$ was also
present, $u_i(t_i, s^*_{-i}) > u_i(s_i, s^*_{-i})$. Repeating from $t_i$ gives a chain with strictly
rising payoffs against $s^*_{-i}$; in a finite game it ends at a never-deleted strategy, $s_i^*$.
So $u_i(s_i^*, s^*_{-i}) > u_i(s_i, s^*_{-i})$ for all $s_i \ne s_i^*$ and all $i$: $s^*$ is a
(strict) Nash equilibrium.}
\end{exsol}
